% Luc Destombe --- licence CC BY-NC-SA 4.0
% https://creativecommons.org/licenses/by-nc-sa/4.0/deed.fr
% Fichier autonome : compiler avec pdflatex, deux fois (signets, nombre de pages).
\def\lyceeversion{1.0}
\documentclass[11pt,a4paper]{article}
\makeatletter
\newif\ifeleve
\elevefalse

\newif\iflycee@cm
\newif\iflycee@palatino

\RequirePackage[utf8]{inputenc}
\RequirePackage[T1]{fontenc}
\RequirePackage[french]{babel}
\RequirePackage{etoolbox}

\newcommand{\ShorthandsOff}{\@ifundefined{shorthandoff}{}{\shorthandoff{;:!?}}}
\newcommand{\ShorthandsOn}{\@ifundefined{shorthandon}{}{\shorthandon{;:!?}}}
\ShorthandsOff
\AtBeginDocument{\ShorthandsOn}
\AtBeginEnvironment{tikzpicture}{\ShorthandsOff}

\RequirePackage{geometry}
\RequirePackage{fancyhdr}
\RequirePackage{lastpage}
\RequirePackage{multicol}
\RequirePackage{array}
\RequirePackage{enumitem}
\RequirePackage{xcolor}
\RequirePackage{needspace}

\RequirePackage{amsmath,amssymb}
\iflycee@cm\else\RequirePackage{fourier}\fi
\DeclareMathAlphabet{\mathcal}{OMS}{cmsy}{m}{n}
\renewcommand{\ttdefault}{lmtt}
\begingroup\catcode`\_=\active \catcode`\^=\active
\global\let_\sb \global\let^\sp
\endgroup
\AtBeginDocument{\catcode`\_=12 \mathcode`\_="8000
  \catcode`\^=12 \mathcode`\^="8000 }
\newlength{\retraitcalculs}
\setlength{\retraitcalculs}{2em}

\RequirePackage{tikz}
\usetikzlibrary{arrows.meta,calc,babel,shapes.misc}
\RequirePackage[most]{tcolorbox}
\tcbuselibrary{skins,breakable}

\definecolor{coulDef}{HTML}{1F4E79}
\definecolor{coulProp}{HTML}{7030A0}
\definecolor{coulRem}{HTML}{C55A11}
\definecolor{coulEx}{HTML}{2E7D32}
\definecolor{coulExo}{HTML}{B03A2E}
\definecolor{coulPart}{HTML}{1F4E79}
\definecolor{coulMeth}{HTML}{1B6E4F}
\definecolor{coulCourbe}{HTML}{0B5ED7}
\definecolor{coulCourbeB}{HTML}{C00000}
\definecolor{coulCourbeC}{HTML}{2E7D32}
\definecolor{coulGrille}{HTML}{8C8C8C}
\definecolor{coulRep}{HTML}{1B6E4F}
\definecolor{coulBar}{HTML}{2E7D32}

\newcommand{\R}{\mathbb{R}}
\newcommand{\Rs}{\mathbb{R}^{*}}
\newcommand{\N}{\mathbb{N}}
\newcommand{\Ns}{\mathbb{N}^{*}}
\newcommand{\Z}{\mathbb{Z}}
\newcommand{\Q}{\mathbb{Q}}
\renewcommand{\le}{\leqslant}
\renewcommand{\ge}{\geqslant}
\renewcommand{\approx}{\simeq}
\newcommand{\vect}[1]{\overrightarrow{#1}}
\newcommand{\norme}[1]{\left\lVert #1 \right\rVert}
\newcommand{\intervalle}[2]{\left[#1\,;#2\right]}
\RequirePackage{cancel}
\renewcommand{\CancelColor}{\color{coulExo}}
\newcommand{\fois}[1]{\times{\color{coulExo}#1}}
\newcommand{\barre}[1]{\cancel{#1}}
\newlength{\lycee@eqh}\newlength{\lycee@eqd}
\newsavebox{\lycee@eqa}\newsavebox{\lycee@eqb}\newsavebox{\lycee@eqc}
\newcommand{\lycee@eqcol}[1]{\hbox{$\renewcommand{\arraystretch}{1.35}%
  \begin{array}[t]{@{}>{\displaystyle}r@{}>{\displaystyle{}}l@{}}#1\end{array}$}}
\newcommand{\lycee@eqmesure}[1]{%
  \ifdim\ht#1>\lycee@eqh\setlength{\lycee@eqh}{\ht#1}\fi
  \ifdim\dp#1>\lycee@eqd\setlength{\lycee@eqd}{\dp#1}\fi}
\newcommand{\lycee@eqtrait}{\hspace{1.2em}\tikz[baseline=0pt]\draw[dash pattern=on 2pt off 2pt,
  line width=0.5pt] (0,\lycee@eqh) -- (0,-\lycee@eqd);\hspace{1.2em}}
\newcommand{\eqproduit}[2]{\par\smallskip\noindent C'est une équation produit :
  $A\times B=0 \iff A=0$ ou $B=0$.\par\nopagebreak\smallskip
  \begingroup\setlength{\lycee@eqh}{0pt}\setlength{\lycee@eqd}{0pt}%
  \setbox\lycee@eqa=\lycee@eqcol{#1}\setbox\lycee@eqb=\lycee@eqcol{#2}%
  \lycee@eqmesure\lycee@eqa\lycee@eqmesure\lycee@eqb%
  \noindent\hspace*{\retraitcalculs}\usebox\lycee@eqa\lycee@eqtrait\usebox\lycee@eqb\endgroup\par\smallskip}
\newcommand{\eqproduittrois}[3]{\par\smallskip\noindent C'est une équation produit :
  $A\times B\times C=0 \iff A=0$ ou $B=0$ ou $C=0$.\par\nopagebreak\smallskip
  \begingroup\setlength{\lycee@eqh}{0pt}\setlength{\lycee@eqd}{0pt}%
  \setbox\lycee@eqa=\lycee@eqcol{#1}\setbox\lycee@eqb=\lycee@eqcol{#2}\setbox\lycee@eqc=\lycee@eqcol{#3}%
  \lycee@eqmesure\lycee@eqa\lycee@eqmesure\lycee@eqb\lycee@eqmesure\lycee@eqc%
  \noindent\hspace*{\retraitcalculs}\usebox\lycee@eqa\lycee@eqtrait\usebox\lycee@eqb\lycee@eqtrait\usebox\lycee@eqc\endgroup
  \par\smallskip}

\tikzset{grille/.style={coulGrille,line width=0.4pt}}
\newcommand{\quadrillage}[4]{\draw[grille,step=1] (#1,#2) grid (#3,#4);}
\tikzset{
  croix/.style={cross out,draw,line width=0.6pt,inner sep=0pt,minimum size=4pt},
  croixdroite/.style={croix,rotate=45,line width=0.8pt},
}
\newcommand{\pt}[3]{\path (#1) node[croix]{} node[#2,font=\small,inner sep=2.5pt]{$#3$};}

\newcommand{\lycee@repere}[4]{%
  \draw[grille,step=1] (#1,#2) grid (#3,#4);
  \draw[-{Stealth[length=2mm]},thick] (#1,0)--({#3+0.4},0) node[below left=-1pt]{$x$};
  \draw[-{Stealth[length=2mm]},thick] (0,#2)--(0,{#4+0.4}) node[below left=-1pt]{$y$};
  \node[below left,font=\scriptsize,inner sep=1.5pt] at (0,0) {$0$};}
\newcommand{\lycee@gradx}[1]{\foreach \k/\l in {#1}{%
  \draw (\k,0.07)--(\k,-0.07) node[below,font=\scriptsize,inner sep=1.5pt]{$\l$};}}
\newcommand{\lycee@grady}[1]{\foreach \k/\l in {#1}{%
  \draw (0.07,\k)--(-0.07,\k) node[left,font=\scriptsize,inner sep=1.5pt]{$\l$};}}
\newcommand{\lycee@intff}[2]{\left[#1\,;#2\right]}
\newcommand{\lycee@intfo}[2]{\left[#1\,;#2\right[}
\newcommand{\lycee@intof}[2]{\left]#1\,;#2\right]}
\newcommand{\lycee@intoo}[2]{\left]#1\,;#2\right[}
\newcommand{\lycee@ens}[1]{\left\{#1\right\}}
\AtBeginDocument{%
  \providecommand{\repere}{\lycee@repere}%
  \providecommand{\gradx}{\lycee@gradx}%
  \providecommand{\grady}{\lycee@grady}%
  \providecommand{\Cf}{\mathcal{C}\sb{f}}%
  \providecommand{\Cg}{\mathcal{C}\sb{g}}%
  \providecommand{\intff}{\lycee@intff}%
  \providecommand{\intfo}{\lycee@intfo}%
  \providecommand{\intof}{\lycee@intof}%
  \providecommand{\intoo}{\lycee@intoo}%
  \providecommand{\ens}{\lycee@ens}}

\newlist{questions}{enumerate}{3}
\setlist[questions,1]{label=\arabic*\textdegree),leftmargin=*,itemsep=2pt,topsep=3pt}
\setlist[questions,2]{label=\alph*),leftmargin=*,itemsep=1pt,topsep=2pt}
\setlist[questions,3]{label=\roman*),leftmargin=*,itemsep=1pt,topsep=1pt}
\setlist[itemize]{label=\textbullet,leftmargin=*,itemsep=2pt,topsep=3pt}

\newcommand{\besoinplace}[1]{\par\penalty\@M\begingroup
  \dimen@=#1\relax\dimen@ii=\pagegoal\advance\dimen@ii-\pagetotal
  \ifdim\dimen@>\dimen@ii\ifdim\dimen@ii>\z@\vfil\fi\break\fi\endgroup}

\newcommand{\lycee@pied}{\thepage/\pageref{LastPage}}
\newcommand{\entete}[2]{%
  \pagestyle{fancy}\fancyhf{}%
  \fancyhead[L]{\footnotesize\color{coulPart}\textbf{#1}}%
  \fancyhead[R]{\footnotesize\color{coulPart}\textbf{#2}}%
  \fancyfoot[C]{\footnotesize\lycee@pied}%
  \renewcommand{\headrulewidth}{0.4pt}%
  \renewcommand{\headrule}{\hbox to\headwidth{%
    \color{coulPart!40}\leaders\hrule height \headrulewidth\hfill}}}

\RequirePackage{ccicons}
\newcommand{\lycee@licence}{{\scriptsize\color{gray}%
  \copyright\ Luc Destombe --- \ccbyncsa\ CC BY-NC-SA 4.0}}
\newif\iflycee@licence \lycee@licencetrue
\newcommand{\sanslicence}{\lycee@licencefalse}
\AtBeginDocument{\iflycee@licence\AddToHookNext{shipout/foreground}{%
  \put(0,0){\rlap{\kern\dimexpr1in+\hoffset+\oddsidemargin+\textwidth\relax
    \llap{\raisebox{-\dimexpr1in+\voffset+\topmargin+\headheight+\headsep
      +\textheight+\footskip\relax}{\lycee@licence}}}}}\fi}

\newcommand{\lycee@site}{faire-des-maths.fr}
\newcommand{\lycee@version}{{\scriptsize\color{gray}\lycee@site
  \ifdefined\lyceeversion\ --- Version~\lyceeversion\fi}}
\AtBeginDocument{\iflycee@licence\AddToHookNext{shipout/foreground}{%
  \put(0,0){\rlap{\kern\dimexpr1in+\hoffset+\oddsidemargin\relax
    \raisebox{-\dimexpr1in+\voffset+\topmargin+\headheight+\headsep
      +\textheight+\footskip\relax}{\lycee@version}}}}\fi}

\newcommand{\formulesserrees}{%
  \setlength{\abovedisplayskip}{3pt}\setlength{\belowdisplayskip}{3pt}%
  \setlength{\abovedisplayshortskip}{2pt}\setlength{\belowdisplayshortskip}{3pt}}

\setlength{\parindent}{0pt}

\RequirePackage[implicit=false,hidelinks,pdfencoding=auto,psdextra,bookmarksopen,
  bookmarksopenlevel=1,pdfcreator={},pdfproducer={}]{hyperref}
\RequirePackage{bookmark}
\hypersetup{pdfauthor={Luc Destombe},pdfkeywords={CC BY-NC-SA 4.0}}
\newcounter{lycee@signet}
\newcommand{\signet}[3][]{\stepcounter{lycee@signet}%
  \edef\lycee@dest{\if\relax\detokenize{#1}\relax
    signet.\arabic{lycee@signet}\else#1\fi}%
  \hypertarget{\lycee@dest}{}\bookmark[level=#2,dest=\lycee@dest]{#3}}
\pdfstringdefDisableCommands{%
  \def\R{R}\def\N{N}\def\Z{Z}\def\Q{Q}\def\vect#1{#1}\def\Cf{Cf}\def\Cg{Cg}\def\fois#1{×#1}%
  \def\barre#1{#1}\def\intervalle#1#2{[#1;#2]}\def\intff#1#2{[#1;#2]}%
  \def\textsuperscript#1{#1}\def\dfrac#1#2{#1/#2}\def\frac#1#2{#1/#2}%
  \def\vec#1{#1⃗}}%
\newcommand{\res}[1]{#1}
\geometry{top=2cm,bottom=1cm,left=1cm,right=1cm,headheight=14pt,footskip=0.6cm}

\RequirePackage{pgfplots}
\pgfplotsset{compat=1.18}
\RequirePackage{tkz-tab}

\newcounter{partiecours}
\renewcommand{\thepartiecours}{\Roman{partiecours}}
\newcommand{\partiecours}[1]{%
  \refstepcounter{partiecours}%
  \Needspace*{8\baselineskip}%
  \bigskip
  \noindent\begin{tcolorbox}[enhanced,colback=coulPart,colframe=coulPart,
      boxrule=0pt,arc=2pt,left=8pt,right=8pt,top=4pt,bottom=4pt,
      before skip=6pt,after skip=10pt]
    \color{white}\bfseries\large \thepartiecours{} --- #1%
    \signet[partie-\arabic{partiecours}]{1}{\thepartiecours{} --- #1}
  \end{tcolorbox}%
  \addcontentsline{toc}{section}{\thepartiecours{} --- #1}%
}

\newtcolorbox{boitecours}[3][]{%
  enhanced,breakable,
  colback=white,colframe=#2,coltitle=white,
  fonttitle=\bfseries,
  attach boxed title to top left={xshift=8pt,yshift=-\tcboxedtitleheight/2},
  boxed title style={colback=#2,arc=2pt,boxrule=0pt},
  boxrule=0.9pt,arc=2pt,
  left=8pt,right=8pt,top=8pt,bottom=6pt,
  title={#3},#1}

\newcounter{methode}
\newenvironment{definitionbox}[1][Définition]%
  {\begin{boitecours}[phantom={\signet{2}{#1}}]{coulDef}{#1}}{\end{boitecours}}
\newenvironment{proprietebox}[1][Propriété]%
  {\begin{boitecours}[phantom={\signet{2}{#1}}]{coulProp}{#1}}{\end{boitecours}}
\newenvironment{methodebox}[1][Méthode]{\stepcounter{methode}%
  \begin{boitecours}[phantom={\signet[methode-\arabic{methode}]{2}{#1}}]{coulMeth}{#1}}%
  {\end{boitecours}}

\newcommand{\etape}[1]{\par\smallskip\textbf{\color{coulMeth}#1}\par\nobreak\smallskip}
\newcommand{\enonce}[1]{\emph{#1}\par\smallskip}
\newcommand{\reponse}[1]{\par\smallskip
  {\footnotesize\itshape\color{black!55}Réponses : #1}}
\newcommand{\demo}[1]{\textbf{\color{coulMeth}Démonstration.} #1}
\newcommand{\afaire}[1]{%
  \par\smallskip
  \noindent{\small\color{coulExo}$\blacktriangleright$\ \textbf{À faire :}
    fiche d'exercices du chapitre, #1.}\par\smallskip}

\newenvironment{calculs}{\setlength{\jot}{5pt}\@fleqntrue\@mathmargin\retraitcalculs
  \setlength{\abovedisplayskip}{4pt}\setlength{\belowdisplayskip}{4pt}%
  \setlength{\abovedisplayshortskip}{4pt}\setlength{\belowdisplayshortskip}{4pt}%
  \csname alignat*\endcsname{2}}%
  {\csname endalignat*\endcsname}
\newcommand{\rem}[1]{\text{\small\itshape\color{black!60}#1}}
\newcommand{\explique}[2]{\par\smallskip\noindent
  \begin{minipage}[t]{0.57\linewidth}\raggedright #1\end{minipage}\hfill
  \begin{minipage}[t]{0.40\linewidth}\raggedright\leavevmode
    {\small\itshape\color{black!60}#2}\end{minipage}\par}
\newcommand{\cas}[1]{\par\smallskip\noindent\textbf{\color{coulExo}#1}\vadjust{\nobreak}}
\newcommand{\calc}[1]{\par\smallskip\textbf{\color{coulExo}#1}\par\nobreak}

\newtcolorbox{remarquebox}[1][Remarque]{%
  enhanced,breakable,colback=white,colframe=coulRem,
  boxrule=0pt,leftrule=2.5pt,arc=0pt,outer arc=0pt,
  left=8pt,right=6pt,top=5pt,bottom=5pt,
  before upper={\textbf{\color{coulRem}#1}\par\smallskip}}

\newtcolorbox{exemplebox}[1][Exemple]{%
  enhanced,breakable,colback=white,colframe=coulEx,
  boxrule=0pt,leftrule=2.5pt,arc=0pt,outer arc=0pt,
  left=8pt,right=6pt,top=5pt,bottom=5pt,
  before upper={\textbf{\color{coulEx}#1}\par\smallskip}}

\pgfplotsset{
  repere/.style={
    axis lines=middle,
    axis line style={-{Stealth[length=2.4mm]},thick},
    every axis x label/.style={at={(ticklabel* cs:1.0)},anchor=west},
    every axis y label/.style={at={(ticklabel* cs:1.0)},anchor=south},
    label style={font=\footnotesize},
    tick label style={font=\scriptsize},
    grid=both,
    major grid style={grille},
    minor tick num=0,
    tick align=inside,
    clip mode=individual,
  },
}
\tikzset{
  courbe/.style={coulCourbe,line width=1.1pt,smooth},
  courbeB/.style={coulCourbeB,line width=1.1pt,smooth},
  courbeC/.style={coulCourbeC,line width=1.1pt,smooth},
}

\newlist{etapes}{enumerate}{1}
\setlist[etapes]{label=\textbf{\arabic*.},leftmargin=*,itemsep=1pt,topsep=2pt}

\newlist{expressions}{enumerate}{1}
\setlist[expressions]{label={\Alph*\ $=$},leftmargin=*,itemsep=3pt,topsep=3pt}

\newcounter{exo}
\renewcommand{\theexo}{\ifnum\value{exo}<10 0\fi\arabic{exo}}
\newtcolorbox{exercice}[1][]{%
  enhanced,breakable,
  colback=white,colframe=coulExo,
  boxrule=0.9pt,arc=2pt,
  left=8pt,right=8pt,top=8pt,bottom=6pt,
  coltitle=white,fonttitle=\bfseries,
  attach boxed title to top left={xshift=8pt,yshift=-\tcboxedtitleheight/2},
  boxed title style={colback=coulExo,arc=2pt,boxrule=0pt},
  phantom={\signet{2}{Exercice \arabic{exo}}},
  title={Exercice \theexo\ifx\relax#1\relax\else{} --- #1\fi}}
\newenvironment{exo}[1][\relax]{\refstepcounter{exo}\begin{exercice}[#1]}{\end{exercice}}

  \newenvironment{correction}{\par}{\par}
  \newcommand{\autableau}{}
  \newcommand{\etapeexemples}{\etape{Exemples corrigés}}
  \newcommand{\etapetableau}[1]{\etape{#1}}
  \newenvironment{exempletableau}[1][Exemple]%
    {\begin{exemplebox}[#1]}{\end{exemplebox}}

\RequirePackage{comment}
\excludecomment{correctionavj}

\newcommand{\coursEntete}{}
\newcommand{\coursNiveau}{}
\newcommand{\coursTitre}{}
\newcommand{\coursSurTitre}{}
\newcommand{\coursSousTitre}{}

\newcommand{\enteteCours}[1]{\renewcommand{\coursEntete}{#1}}
\newcommand{\chapitrecours}[2]{\enteteCours{Chapitre #1 --- #2}}
\newcommand{\niveaucours}[1]{\renewcommand{\coursNiveau}{#1}}
\newcommand{\titrecours}[3][]{%
  \renewcommand{\coursTitre}{#2}%
  \renewcommand{\coursSurTitre}{#1}%
  \renewcommand{\coursSousTitre}{#3}}

\pagestyle{fancy}
\fancyhf{}
\fancyhead[L]{\footnotesize\color{coulPart}\textbf{\coursEntete}}
\fancyhead[R]{\footnotesize\color{coulPart}\textbf{\coursNiveau}}
\fancyfoot[C]{\footnotesize\thepage}
\renewcommand{\headrulewidth}{0.4pt}
\renewcommand{\headrule}{\hbox to\headwidth{\color{coulPart!40}\leaders\hrule height \headrulewidth\hfill}}

\setlength{\parskip}{2pt}

\newcommand{\lycee@mentioneleve}{\\[2pt]{\small\itshape Version élève}}
\AtBeginDocument{%
  \begin{center}
    \begin{tcolorbox}[enhanced,width=0.72\linewidth,colback=white,colframe=coulPart,
        boxrule=1.6pt,arc=3pt,halign=center,top=10pt,bottom=10pt]
      \ifx\coursSurTitre\empty
        {\Huge\bfseries\color{coulPart} \coursTitre}\\[4pt]
      \else
        {\Huge\bfseries\color{coulPart} \coursTitre}\\[3pt]
        {\Large\bfseries\color{coulPart} \coursSurTitre}\\[5pt]
      \fi
      {\normalsize \coursSousTitre}%
      
    \end{tcolorbox}
  \end{center}%
}

\RequirePackage{tikz}
\RequirePackage{xcolor}

\usetikzlibrary{decorations.text}
\definecolor{lycee@metal}{HTML}{1F4E79}
\definecolor{lycee@manche}{HTML}{8A5A2B}
\definecolor{lycee@etincelle}{HTML}{F28C28}
\newcommand{\lycee@pluriel}[2]{#1 #2\ifnum#1>1 s\fi}
\newcommand{\lycee@sceau}{%
  \edef\lycee@bas{\lycee@pluriel{\lyceerelectures}{relecture} et
    \lycee@pluriel{\lyceerodages}{rodage}}%
  \begin{tikzpicture}[line join=round]
    \fill[coulPart!10,draw=coulPart,line width=.8pt] (0:2.2)
      \foreach \i in {1,...,71}{ -- (\i*5:{isodd(\i)?2.1:2.2}) } -- cycle;
    \draw[coulPart,line width=1.2pt] (0,0) circle (1.93);
    \draw[coulPart,line width=.4pt] (0,0) circle (1.85);
    \draw[coulPart,line width=.4pt] (0,0) circle (1.2);
    \path[decorate,decoration={text along path,text align=center,
        text={|\fontsize{9.5}{11}\selectfont\bfseries\color{coulPart}|TEST{\'E} AVEC MES {\'E}L{\`E}VES}}]
      (180:1.45) arc (180:0:1.45);
    \path[decorate,decoration={text along path,text align=center,
        text/.expanded={|\noexpand\fontsize{8.5}{10}\noexpand\selectfont\noexpand\color{coulPart}|\lycee@bas}}]
      (195:1.68) arc (195:345:1.68);
    \begin{scope}[x=.25mm,y=-.25mm]
      \begin{scope}[shift={(-32,-40)}]%
        \fill[lycee@metal] (23,31) rectangle (41,49);
        \draw[lycee@etincelle,line width=.5pt,line cap=round]
          (20,30)--(15,27) (19,34)--(14,35) (44,30)--(49,27) (45,34)--(50,35);
        \fill[lycee@manche,rounded corners=.3mm] (38,24.6) rectangle (60,27.6);
        \fill[lycee@metal,rounded corners=.4mm] (25,20) rectangle (39,31);
        \fill[lycee@etincelle,rounded corners=.2mm] (25,20) rectangle (28,31);
      \end{scope}
    \end{scope}
    \ifdefined\lyceesceauplusieurs\def\lycee@annee{Ann{\'e}es~:}\else
      \def\lycee@annee{Ann{\'e}e}\fi
    \path[decorate,decoration={text along path,text align=center,
        text/.expanded={|\noexpand\fontsize{8}{10}\noexpand\selectfont\noexpand\color{coulPart}|%
          \unexpanded\expandafter{\lycee@annee} \unexpanded\expandafter{\lyceesceauannees}}}]
      (95:1.08) arc (95:445:1.08);
  \end{tikzpicture}}
\AtEndDocument{\ifdefined\lyceesceauannees
  \par\besoinplace{5cm}\vspace{1.5em}%
  \begin{center}\lycee@sceau\end{center}\fi}
\makeatother

\chapitrecours{3}{Polynômes du second degré}
\niveaucours{1\textsuperscript{re} STMG}
\titrecours{POLYNÔMES DU SECOND DEGRÉ}{Cours --- Première STMG}

\begin{document}

\begin{remarquebox}[Comment utiliser ce cours]
Chaque partie commence par ce qu'il faut \textbf{savoir} (définitions, propriétés), puis
vient ce qu'il faut \textbf{savoir faire} : une méthode, avec des
\textbf{exemples corrigés} faits en classe, puis des exercices
\og\textbf{À vous de jouer}\fg{} à chercher seul. Tout le chapitre se fait
\textbf{sans calculatrice}.

\end{remarquebox}

\partiecours{Fonction polynôme du second degré}

\begin{definitionbox}[Définition 1 --- Fonction polynôme du second degré]
Une \textbf{fonction polynôme du second degré} est une fonction $f$ de la forme
\[
  f(x)=ax^{2}+bx+c \qquad\text{avec } a\ne 0 .
\]
Les nombres $a$, $b$ et $c$ sont ses \textbf{coefficients}. Sa courbe s'appelle une
\textbf{parabole}.
\end{definitionbox}

\begin{remarquebox}
Sans terme en $x^{2}$, il reste $bx+c$ : c'est une fonction affine, dont la courbe est
une droite.
\end{remarquebox}

\begin{methodebox}[Méthode 1 --- Reconnaître une fonction du second degré, calculer une image]
\etapeexemples
\begin{questions}[label=\alph*)]
  \item Pour chaque fonction $f$, $g$ et $h$, dire si elle est du second degré et, si
        oui, donner $a$, $b$ et $c$ :
        $f(x)=3x^{2}-2x+1$ \; ; \; $g(x)=5-x^{2}$ \; ; \; $h(x)=4x+7$.
  \item Calculer $f(-2)$.
  \item Les points $A(1\,;2)$ et $B(-1\,;4)$ sont-ils sur la courbe de $f$ ?
\end{questions}
\begin{correction}
\cas{a)}
\explique{$f$ : \res{$a=3$ ; $b=-2$ ; $c=1$}.}{chaque coefficient \textbf{avec son signe}}
\explique{$g(x)=-x^{2}+5$ : \res{$a=-1$ ; $b=0$ ; $c=5$}.}{on range ; pas de terme en $x$ :
$b=0$}
\explique{$h$ \res{n'est pas du second degré} : c'est une fonction affine.}{pas de $x^{2}$}
\cas{b)}
\begin{calculs}
f(-2) &= 3\times(-2)^{2}-2\times(-2)+1 &\qquad& \rem{on remplace $x$ par $(-2)$}\\
f(-2) &= 3\times 4+4+1 &\qquad& \rem{$(-2)^{2}=4$}\\
f(-2) &= \res{17} &&
\end{calculs}
\cas{c)}
\explique{$f(1)=3-2+1=2$ : c'est l'ordonnée de $A$, donc \res{$A$ est sur la
courbe}.}{on calcule l'image de l'abscisse}
\explique{$f(-1)=3+2+1=6$ et $6\ne 4$ : \res{$B$ n'est pas sur la courbe}.}{}
\end{correction}

\etape{À vous de jouer}
\begin{questions}[label=\alph*)]
  \item Pour chaque fonction $u$, $v$ et $w$, dire si elle est du second degré et, si
        oui, donner $a$, $b$ et $c$ :
        $u(x)=-2x^{2}+x+4$ \; ; \; $v(x)=x^{2}-3x$ \; ; \; $w(x)=1-6x$.
  \item Calculer $u(-1)$.
  \item Les points $C(2\,;-2)$ et $D(-2\,;0)$ sont-ils sur la courbe de $u$ ?
\end{questions}
\reponse{a) $u$ : $a=-2$ ; $b=1$ ; $c=4$ ;\; $v$ : $a=1$ ; $b=-3$ ; $c=0$ ;\; $w$ :
affine ;\; b)~$u(-1)=1$ ;\; c)~$C$ oui ; $D$ non ($u(-2)=-6$).}
\end{methodebox}

\afaire{exercices 1 à 4}

\partiecours{Forme factorisée et racines}

\begin{definitionbox}[Définition 2 --- Forme factorisée, racines]
\begin{minipage}[c]{0.58\linewidth}\raggedright
Une fonction du second degré est écrite sous \textbf{forme factorisée} quand elle
s'écrit $f(x)=a(x-x_{1})(x-x_{2})$.

\smallskip
Les \textbf{racines} de $f$ sont les solutions de $f(x)=0$. Ce sont les abscisses des
points où la parabole coupe l'axe des abscisses.
\end{minipage}\hfill
\begin{minipage}[c]{0.38\linewidth}\centering
\begin{tikzpicture}[x=0.42cm,y=0.42cm]
  \repere{-3}{-5}{5}{4}
  \gradx{-2,-1,1,2,3,4} \grady{-4,-2,2}
  \draw[coulCourbe,line width=1.1pt,domain=-1.9:3.9,samples=60] plot (\x,{(\x+1)*(\x-3)});
  \path[red] (-1,0) node[croix]{} (3,0) node[croix]{};
  \node[red,above left,font=\footnotesize,inner sep=1pt] at (-1,0) {$x_{1}$};
  \node[red,above right,font=\footnotesize,inner sep=1pt] at (3,0) {$x_{2}$};
\end{tikzpicture}
\end{minipage}
\end{definitionbox}

\begin{proprietebox}[Propriété 1 --- Racines d'une forme factorisée]
Les racines de $f(x)=a(x-x_{1})(x-x_{2})$ sont $x_{1}$ et $x_{2}$ : chacune annule une
parenthèse. Dans $(x-4)$, la racine est $4$ ; dans $(x+2)$, elle est $-2$.
\end{proprietebox}

\begin{methodebox}[Méthode 2 --- Utiliser une forme factorisée]
\etapeexemples
Soit $f(x)=2x^{2}+4x-6$.
\begin{questions}[label=\alph*)]
  \item Vérifier que $f(x)=2(x-1)(x+3)$.
  \item En déduire les racines de $f$.
  \item Donner les points où la courbe de $f$ coupe l'axe des abscisses.
\end{questions}
\begin{correction}
\cas{a)}
\begin{calculs}
2(x-1)(x+3) &= 2\left(x^{2}+3x-x-3\right) &\qquad& \rem{on développe la forme
factorisée : d'abord les parenthèses}\\
2(x-1)(x+3) &= 2\left(x^{2}+2x-3\right) &&\\
2(x-1)(x+3) &= 2x^{2}+4x-6 &\qquad& \rem{puis on multiplie par $2$}
\end{calculs}
\explique{On retrouve $f(x)$ : donc $f(x)=2(x-1)(x+3)$.}{« vérifier » : on part de
l'expression la plus compliquée}
\cas{b)}
\begin{calculs}
2(x-1)(x+3) &= 0 &\qquad& \rem{on résout $f(x)=0$ ; le $2$ n'est jamais nul}
\end{calculs}
\eqproduit{x-1 &= 0\\ x &= 1}{x+3 &= 0\\ x &= -3}
\explique{Les racines de $f$ sont \res{$-3$ et $1$}.}{}
\cas{c)}
\explique{La courbe coupe l'axe des abscisses aux points \res{$(-3\,;0)$} et
\res{$(1\,;0)$}.}{ordonnée nulle sur l'axe des abscisses}
\end{correction}

\etape{À vous de jouer}
Mêmes questions pour :
\begin{questions}[label=\alph*)]
  \item $g(x)=-x^{2}+2x+8$, avec $g(x)=-(x-4)(x+2)$ ;
  \item $h(x)=3x^{2}-12$, avec $h(x)=3(x-2)(x+2)$.
\end{questions}
\reponse{a) racines $-2$ et $4$ ; points $(-2\,;0)$ et $(4\,;0)$ ;\; b) racines $-2$ et
$2$ ; points $(-2\,;0)$ et $(2\,;0)$.}
\end{methodebox}

\afaire{exercices 5 à 7}

\partiecours{Représentation graphique}

\begin{proprietebox}[Propriété 2 --- Allure de la parabole]
La parabole de $f(x)=ax^{2}+bx+c$ :
\begin{itemize}
  \item est tournée vers le \textbf{haut} si $a>0$, vers le \textbf{bas} si $a<0$ ;
  \item a un \textbf{axe de symétrie} vertical, qui passe par son point le plus bas
        (ou le plus haut) : le \textbf{sommet} $S$ ;
  \item coupe l'axe des ordonnées au point $(0\,;c)$, car $f(0)=c$.
\end{itemize}
\begin{center}
\begin{minipage}[t]{0.47\linewidth}\centering
{\small $a>0$ : $f(x)=x^{2}-2x-1$}\par\smallskip
\begin{tikzpicture}[x=0.42cm,y=0.42cm]
  \repere{-2}{-3}{4}{4}
  \gradx{-1,1,2,3} \grady{-2,-1,1,2,3}
  \draw[coulCourbe,line width=1.1pt,domain=-1.4:3.4,samples=60] plot (\x,{\x*\x-2*\x-1});
  \draw[red,dashed,thick] (1,-3)--(1,4);
  \path (1,-2) node[croix]{} node[below right,font=\footnotesize,inner sep=2pt]{$S$};
  \path (0,-1) node[croix]{};
\end{tikzpicture}
\end{minipage}\hfill
\begin{minipage}[t]{0.47\linewidth}\centering
{\small $a<0$ : $g(x)=-x^{2}+2x+2$}\par\smallskip
\begin{tikzpicture}[x=0.42cm,y=0.42cm]
  \repere{-2}{-3}{4}{4}
  \gradx{-1,1,2,3} \grady{-2,-1,1,2,3}
  \draw[coulCourbe,line width=1.1pt,domain=-1.2:3.2,samples=60] plot (\x,{-\x*\x+2*\x+2});
  \draw[red,dashed,thick] (1,-3)--(1,4);
  \path (1,3) node[croix]{} node[above right,font=\footnotesize,inner sep=2pt]{$S$};
  \path (0,2) node[croix]{};
\end{tikzpicture}
\end{minipage}
\end{center}
\end{proprietebox}

\begin{methodebox}[Méthode 3 --- Tracer une parabole]
\etapeexemples
Soit $f(x)=x^{2}-2x-3$.
\begin{questions}[label=\alph*)]
  \item Compléter le tableau de valeurs de $f$ pour les entiers de $-2$ à $4$.
  \item Tracer la courbe de $f$.
  \item Lire le sommet et l'axe de symétrie.
\end{questions}
\begin{correction}
\cas{a)}\par\nopagebreak
\begin{minipage}[c]{0.58\linewidth}\raggedright
\explique{$f(-2)=4+4-3=5$ ; $f(-1)=1+2-3=0$ ; …}{on remplace $x$ par chaque valeur}
\begin{center}\small
\renewcommand{\arraystretch}{1.3}
\begin{tabular}{|c|c|c|c|c|c|c|c|}
\hline
$x$ & $-2$ & $-1$ & $0$ & $1$ & $2$ & $3$ & $4$\\
\hline
$f(x)$ & $5$ & $0$ & $-3$ & $-4$ & $-3$ & $0$ & $5$\\
\hline
\end{tabular}
\end{center}
\cas{b)}
\explique{On place les points, puis on les relie par une courbe arrondie, sans
règle.}{$a=1>0$ : tournée vers le haut}
\cas{c)}
\explique{Sommet \res{$S(1\,;-4)$} ; axe de symétrie : \res{la droite verticale qui passe par
$x=1$}.}{les valeurs se répètent de part et d'autre de $1$}
\end{minipage}\hfill
\begin{minipage}[c]{0.38\linewidth}\centering
\begin{tikzpicture}[x=0.4cm,y=0.4cm]
  \repere{-3}{-5}{5}{6}
  \gradx{-2,-1,1,2,3,4} \grady{-4,-2,2,4}
  \draw[coulCourbe,line width=1.1pt,domain=-2:4,samples=60] plot (\x,{\x*\x-2*\x-3});
  \foreach \p in {(-2,5),(-1,0),(0,-3),(1,-4),(2,-3),(3,0),(4,5)}
    \path[coulCourbe] \p node[croix]{};
  \draw[red,dashed,thick] (1,-5)--(1,6);
  \node[below right,font=\footnotesize,inner sep=2pt] at (1,-4) {$S$};
\end{tikzpicture}
\end{minipage}
\end{correction}

\etape{À vous de jouer}
Mêmes questions pour $g(x)=-x^{2}+4x$, avec les entiers de $-1$ à $5$.
\reponse{$g(-1)=-5$ ; $g(0)=0$ ; $g(1)=3$ ; $g(2)=4$ ; $g(3)=3$ ; $g(4)=0$ ; $g(5)=-5$ ;\;
sommet $S(2\,;4)$, axe : la droite verticale qui passe par $x=2$.}
\end{methodebox}

\begin{methodebox}[Méthode 4 --- Associer une fonction à sa parabole]
\etapeexemples
\begin{minipage}[c]{0.58\linewidth}\raggedright
Associer chaque fonction à sa courbe :
\[
  f(x)=x^{2}+1 \qquad g(x)=-x^{2}+3 \qquad h(x)=x^{2}-2
\]
\end{minipage}\hfill
\begin{minipage}[c]{0.38\linewidth}\centering
\begin{tikzpicture}[x=0.42cm,y=0.42cm]
  \repere{-3}{-3}{3}{5}
  \gradx{-2,2} \grady{-2,2,4}
  \draw[coulCourbe,line width=1.1pt,domain=-2.6:2.6,samples=60] plot (\x,{\x*\x-2});
  \draw[coulCourbeB,line width=1.1pt,domain=-2.2:2.2,samples=60] plot (\x,{-\x*\x+3});
  \draw[coulCourbeC,line width=1.1pt,domain=-2:2,samples=60] plot (\x,{\x*\x+1});
  \node[coulCourbe,font=\footnotesize] at (2.9,3.6) {$\mathcal{C}_{1}$};
  \node[coulCourbeB,font=\footnotesize] at (2.6,-2.4) {$\mathcal{C}_{2}$};
  \node[coulCourbeC,font=\footnotesize] at (-1.2,4.6) {$\mathcal{C}_{3}$};
\end{tikzpicture}
\end{minipage}
\begin{correction}
\explique{$g$ est la seule avec $a<0$ : sa courbe est tournée vers le bas, c'est
\res{$\mathcal{C}_{2}$}.}{le signe de $a$}
\explique{$f(0)=1$ et $h(0)=-2$ : $f$ correspond à \res{$\mathcal{C}_{3}$}, qui passe par
$(0\,;1)$, et $h$ à \res{$\mathcal{C}_{1}$}, qui passe par $(0\,;-2)$.}{le point $(0\,;c)$ sur
l'axe des ordonnées}
\end{correction}

\etape{À vous de jouer}
\begin{minipage}[c]{0.58\linewidth}\raggedright
Associer chaque fonction à sa courbe :
\[
  p(x)=-x^{2}-1 \qquad q(x)=x^{2}-3 \qquad r(x)=-x^{2}+2
\]
\reponse{$p$ : $\mathcal{C}_{3}$ ;\; $q$ : $\mathcal{C}_{2}$ ;\; $r$ : $\mathcal{C}_{1}$.}
\end{minipage}\hfill
\begin{minipage}[c]{0.38\linewidth}\centering
\begin{tikzpicture}[x=0.42cm,y=0.42cm]
  \repere{-3}{-5}{3}{4}
  \gradx{-2,2} \grady{-4,-2,2}
  \draw[coulCourbe,line width=1.1pt,domain=-2.6:2.6,samples=60] plot (\x,{-\x*\x+2});
  \draw[coulCourbeB,line width=1.1pt,domain=-2.6:2.6,samples=60] plot (\x,{\x*\x-3});
  \draw[coulCourbeC,line width=1.1pt,domain=-2:2,samples=60] plot (\x,{-\x*\x-1});
  \node[coulCourbe,font=\footnotesize] at (-1.2,3) {$\mathcal{C}_{1}$};
  \node[coulCourbeB,font=\footnotesize] at (2.9,2.6) {$\mathcal{C}_{2}$};
  \node[coulCourbeC,font=\footnotesize] at (2.6,-4.4) {$\mathcal{C}_{3}$};
\end{tikzpicture}
\end{minipage}
\end{methodebox}

\afaire{exercices 8 à 10}

\partiecours{Tableaux de variations}

\begin{proprietebox}[Propriété 3 --- Sommet et variations]
\begin{minipage}[c]{0.58\linewidth}\raggedright
Si $f$ a deux racines $x_{1}$ et $x_{2}$, le sommet $S$ de la parabole est \textbf{au
milieu} : son abscisse est
\[
  x_{S}=\dfrac{x_{1}+x_{2}}{2}
\]
et son ordonnée est $f(x_{S})$.
\end{minipage}\hfill
\begin{minipage}[c]{0.38\linewidth}\centering
\begin{tikzpicture}[x=0.42cm,y=0.42cm]
  \repere{-3}{-5}{5}{3}
  \gradx{-2,-1,1,2,3,4} \grady{-4,-2,2}
  \draw[coulCourbe,line width=1.1pt,domain=-1.9:3.9,samples=60] plot (\x,{(\x+1)*(\x-3)});
  \draw[red,dashed,thick] (1,-5)--(1,3);
  \path[red] (-1,0) node[croix]{} (3,0) node[croix]{};
  \node[red,above left,font=\footnotesize,inner sep=1pt] at (-1,0) {$x_{1}$};
  \node[red,above right,font=\footnotesize,inner sep=1pt] at (3,0) {$x_{2}$};
  \path (1,-4) node[croix]{} node[below right,font=\footnotesize,inner sep=2pt]{$S$};
\end{tikzpicture}
\end{minipage}

\smallskip
\begin{minipage}[t]{0.47\linewidth}\centering
{\small $a>0$ : $f(x_{S})$ est le \textbf{minimum}}\par\smallskip
\begin{tikzpicture}[scale=0.85]
  \tkzTabInit[lgt=1.6,espcl=2.2]{$x$/1,$f(x)$/1.6}{,$x_{S}$,}
  \tkzTabVar{+/ ,-/$f(x_{S})$,+/ }
\end{tikzpicture}
\end{minipage}\hfill
\begin{minipage}[t]{0.47\linewidth}\centering
{\small $a<0$ : $f(x_{S})$ est le \textbf{maximum}}\par\smallskip
\begin{tikzpicture}[scale=0.85]
  \tkzTabInit[lgt=1.6,espcl=2.2]{$x$/1,$f(x)$/1.6}{,$x_{S}$,}
  \tkzTabVar{-/ ,+/$f(x_{S})$,-/ }
\end{tikzpicture}
\end{minipage}
\end{proprietebox}

\begin{methodebox}[Méthode 5 --- Dresser le tableau de variations]
\etapeexemples
Donner le sommet de la parabole, puis le tableau de variations de :
\begin{questions}[label=\alph*)]
  \item $f(x)=2(x-1)(x+3)$
  \item $g(x)=-(x-2)(x-6)$
  \item $h(x)=3x^{2}-12x$
\end{questions}
\begin{correction}
\cas{a)}
\explique{Racines : $1$ et $-3$.}{on annule chaque parenthèse}
\begin{calculs}
x_{S} &= \dfrac{1+(-3)}{2} &\qquad& \rem{le milieu des racines}\\
x_{S} &= -1 &&
\end{calculs}
\explique{$f(-1)=2\times(-2)\times 2=-8$ : sommet \res{$S(-1\,;-8)$}.}{on remplace $x$ par $-1$
dans la forme factorisée}
\explique{$a=2>0$ : $f$ diminue, puis augmente ; son minimum est $-8$.}{tournée vers le
haut}
\begin{center}
\begin{tikzpicture}[scale=0.85]
  \tkzTabInit[lgt=1.6,espcl=2.2]{$x$/1,$f(x)$/1.6}{,$-1$,}
  \tkzTabVar{+/ ,-/$-8$,+/ }
\end{tikzpicture}
\end{center}
\cas{b)}
\explique{Racines : $2$ et $6$ ; $x_{S}=\dfrac{2+6}{2}=4$.}{le milieu des racines}
\explique{$g(4)=-(2)\times(-2)=4$ : sommet \res{$S(4\,;4)$}.}{$4-2=2$ et $4-6=-2$}
\explique{$a=-1<0$ : $g$ augmente, puis diminue ; son maximum est $4$.}{tournée vers le
bas}
\begin{center}
\begin{tikzpicture}[scale=0.85]
  \tkzTabInit[lgt=1.6,espcl=2.2]{$x$/1,$g(x)$/1.6}{,$4$,}
  \tkzTabVar{-/ ,+/$4$,-/ }
\end{tikzpicture}
\end{center}
\cas{c)}
\begin{calculs}
3x^{2}-12x &= 0 &\qquad& \rem{les racines : on résout $h(x)=0$}\\
x(3x-12) &= 0 &\qquad& \rem{on factorise par $x$}
\end{calculs}
\eqproduit{x &= 0}{3x-12 &= 0\\ 3x &= 12\\ x &= 4}
\explique{$x_{S}=\dfrac{0+4}{2}=2$ et $h(2)=12-24=-12$ : sommet
\res{$S(2\,;-12)$}.}{$3\times 2^{2}=12$}
\explique{$a=3>0$ : $h$ diminue, puis augmente ; son minimum est $-12$.}{}
\begin{center}
\begin{tikzpicture}[scale=0.85]
  \tkzTabInit[lgt=1.6,espcl=2.2]{$x$/1,$h(x)$/1.6}{,$2$,}
  \tkzTabVar{+/ ,-/$-12$,+/ }
\end{tikzpicture}
\end{center}
\end{correction}

\etape{À vous de jouer}
Donner le sommet de la parabole, puis le tableau de variations de :
\[
  f(x)=3(x-2)(x+4) \qquad g(x)=-2(x-1)(x-5) \qquad h(x)=-x^{2}+8x
\]
\reponse{$f$ : $S(-1\,;-27)$, minimum ;\; $g$ : $S(3\,;8)$, maximum ;\; $h$ : $S(4\,;16)$,
maximum.}
\end{methodebox}

\afaire{exercices 11 à 13}

\partiecours{Problèmes type épreuve anticipée}

À l'épreuve, une fonction du second degré sert de \textbf{modèle} (bénéfice, hauteur,
population…) : on calcule des images, on \textbf{vérifie} une forme factorisée donnée,
puis on en tire les racines, le sommet et les variations, et on \textbf{interprète}
chaque résultat dans le contexte.

\begin{methodebox}[Méthode 6 --- Étudier un modèle du second degré]
\etapeexemples
Un atelier fabrique et vend entre $0$ et $1\,200$ trottinettes électriques par an. On
note $x$ le nombre de trottinettes, en centaines ($x$ entre $0$ et $12$). Le
bénéfice, en milliers d'euros, est modélisé par
\[
  B(x)=-2x^{2}+24x-40 .
\]
Sa courbe est tracée ci-dessous ($x$ en centaines, $B(x)$ en milliers d'euros).
\begin{center}
\begin{tikzpicture}[x=0.7cm,y=0.4cm]
  \repere{0}{-4}{13}{5}
  \gradx{1,2,3,4,5,6,7,8,9,10,11,12}
  \grady{-4/-40,-3/-30,-2/-20,-1/-10,1/10,2/20,3/30,4/40}
  \draw[coulCourbe,line width=1.1pt,domain=0:12,samples=80] plot (\x,{(-2*\x*\x+24*\x-40)/10});
\end{tikzpicture}
\end{center}
\begin{questions}
  \item Par lecture graphique, déterminer pour quelles quantités le bénéfice est de
        $30\,000$~€.
  \item Calculer $B(1)$. Interpréter ce résultat dans le contexte de l'exercice.
  \item Vérifier que $B(x)=-2(x-2)(x-10)$.
  \item Pour quelles quantités le bénéfice est-il nul ?
  \item Dresser le tableau de variations de $B$ entre $0$ et $12$.
  \item Quel est le bénéfice maximal ? Pour combien de trottinettes ?
\end{questions}
\begin{correction}
\cas{1\textdegree)}
\explique{La droite d'ordonnée $30$ coupe la courbe en $x=5$ et $x=7$.}{$30$ milliers
d'euros : la graduation $30$}
\explique{Le bénéfice est de $30\,000$~€ pour \res{$500$ ou $700$ trottinettes}.}{on revient aux
unités de l'énoncé}
\cas{2\textdegree)}
\begin{calculs}
B(1) &= -2\times 1^{2}+24\times 1-40 &\qquad& \rem{on remplace $x$ par $1$}\\
B(1) &= \res{-18} &&
\end{calculs}
\explique{Pour $100$ trottinettes, l'atelier \res{perd $18\,000$~€}.}{$x=1$ : une centaine ;
$B(1)<0$ : une perte}
\cas{3\textdegree)}
\begin{calculs}
-2(x-2)(x-10) &= -2\left(x^{2}-10x-2x+20\right) &\qquad& \rem{on développe la forme
donnée}\\
-2(x-2)(x-10) &= -2\left(x^{2}-12x+20\right) &&\\
-2(x-2)(x-10) &= -2x^{2}+24x-40 &\qquad& \rem{on retrouve $B(x)$}
\end{calculs}
\cas{4\textdegree)}
\explique{On résout $-2(x-2)(x-10)=0$.}{la forme factorisée}
\eqproduit{x-2 &= 0\\ x &= 2}{x-10 &= 0\\ x &= 10}
\explique{Le bénéfice est nul pour \res{$200$ ou $1\,000$ trottinettes}.}{}
\cas{5\textdegree)}
\explique{$x_{S}=\dfrac{2+10}{2}=6$ et $B(6)=-2\times 4\times(-4)=32$.}{le milieu des
racines}
\explique{$B(0)=-40$ et $B(12)=-2\times 10\times 2=-40$.}{les valeurs aux bouts}
\explique{$a=-2<0$ : $B$ \res{augmente de $0$ à $6$, puis diminue de $6$ à $12$}.}{tournée
vers le bas}
\begin{center}
\begin{tikzpicture}[scale=0.85]
  \tkzTabInit[lgt=1.6,espcl=2.2]{$x$/0.8,$B(x)$/1.4}{$0$,$6$,$12$}
  \tkzTabVar{-/$-40$,+/$32$,-/$-40$}
\end{tikzpicture}
\end{center}
\cas{6\textdegree)}
\explique{Le bénéfice maximal est de \res{$32\,000$~€}, pour \res{$600$ trottinettes}.}{on revient
aux unités de l'énoncé}
\end{correction}

\etape{À vous de jouer}
\begin{minipage}[c]{0.58\linewidth}\raggedright
On lance une fusée à eau depuis le sol. Sa hauteur, en mètres, au bout de $t$ secondes
($t$ entre $0$ et $4$) est
\[
  h(t)=-5t^{2}+20t .
\]
Sa courbe est tracée ci-contre ($t$ en secondes, $h(t)$ en mètres).
\end{minipage}\hfill
\begin{minipage}[c]{0.38\linewidth}\centering
\begin{tikzpicture}[x=1cm,y=0.45cm]
  \draw[grille,step=1] (0,0) grid (5,5);
  \draw[-{Stealth[length=2mm]},thick] (0,0)--(5.4,0) node[below left=-1pt]{$t$};
  \draw[-{Stealth[length=2mm]},thick] (0,0)--(0,5.4) node[below left=-1pt]{$h$};
  \node[below left,font=\scriptsize,inner sep=1.5pt] at (0,0) {$0$};
  \gradx{1,2,3,4}
  \grady{1/5,2/10,3/15,4/20}
  \draw[coulCourbe,line width=1.1pt,domain=0:4,samples=60] plot (\x,{(-5*\x*\x+20*\x)/5});
\end{tikzpicture}
\end{minipage}
\begin{questions}
  \item Par lecture graphique, déterminer à quels instants la fusée est à $15$~m de
        haut.
  \item Calculer $h(1)$. Interpréter ce résultat dans le contexte de l'exercice.
  \item Vérifier que $h(t)=-5t(t-4)$.
  \item À quels instants la fusée est-elle au sol ?
  \item Dresser le tableau de variations de $h$ entre $0$ et $4$.
  \item Quelle est la hauteur maximale ? À quel instant ?
\end{questions}
\reponse{1\textdegree) au bout de $1$~s et de $3$~s ;\; 2\textdegree) $h(1)=15$ : $15$~m au
bout de $1$~s ;\; 4\textdegree) $t=0$ (le départ) et $t=4$ (la retombée) ;\;
5\textdegree) croissante de $0$ à $2$, décroissante de $2$ à $4$ ;\; 6\textdegree) $20$~m
au bout de $2$~s.}
\end{methodebox}

\afaire{exercices 14 à 16, puis les exercices de type épreuve 17 à 19}

\end{document}
